\section*{Verzeichnis der KI-Nutzung}
\addcontentsline{toc}{section}{Verzeichnis der KI-Nutzung}
\thispagestyle{plain}

Gemäß den Vorgaben des Seminars \textit{Gesundheitstechnik} (SoSe 2026) wird die Nutzung
generativer KI-Werkzeuge nachfolgend dokumentiert.

\begin{table}[htbp]
    \centering
    \caption{KI-Nutzung im vorliegenden Fachartikel}
    \label{tab:ki_nutzung}
    \begin{tabularx}{\textwidth}{|l|l|X|}
        \hline
        \textbf{Abschnitt} & \textbf{Tool} & \textbf{Zweck und Prompt (sinngemäß)} \\
        \hline
        Literaturrecherche &
        Claude claude-sonnet-4-6 &
        Suche nach Primärquellen zu KI-Bildanalyse, Federated Learning und
        Datensouveränität im Gesundheitswesen (05.05.2026). Prompt: \textit{„Recherchiere
        akademische Quellen zu KI-Bildanalyse in Industrie und Medizin, Modelltraining,
        Benchmarks und Inhouse-Architekturen."} \\
        \hline
        Textentwurf &
        Claude claude-sonnet-4-6 &
        Erstellung des initialen Textentwurfs auf Basis vorbereiteter
        Rechercheergebnisse und Gliederungsvorgabe. Prompt: \textit{„Schreibe einen
        wissenschaftlichen Fachartikel auf ca. 2 Seiten zu KI-Bildanalyse in Industrie
        und Medizin -- Modelltraining und Messbarkeit -- sachlich, ohne Ich-Form,
        IEEE-Zitierweise."} Alle Inhalte wurden vom Autor geprüft. \\
        \hline
    \end{tabularx}
\end{table}

\textit{Hinweis:} Der Autor trägt die volle Verantwortung für die inhaltliche Richtigkeit
aller Aussagen.
